% Pursuing Stacks (cotes 134-2 to 134-8): archive cote 134-8, site page 11 (PDF page 12), typescript page 565.
% typescript. Draft compared: ps7.tex:504-567.
% AI-transcribed by Claude Opus 5.5 (claude-opus-5-5, adaptive thinking, effort high), Message Batches API, batch ps4 (msgbatch_012Aoc1NuskATPsyBeKSYK7J), 2026-11-12; not yet checked by eye.
\scanpage{134-8/11}
\mspage{565}

27.10.

136.    I just spent a couple of hours,  after reading the notes of last
night,  trying to get a better feeling of the basic homology operation
in the context of the basic localizer (Cat),  namely taking the left
derived functor  $Lf_!^{ab}$ of the ``unusual'' direct image functor for abelian
presheaves,  associated to a map
\[ f: A \longrightarrow B \]
in (Cat).  This led me to read again the notes of section 92,  when I
unsuspectingly started on an ``afterthought''\hand{,}\inserted{later gradually} turning into a systematic
reflection on abelianization.  With a distance of nearly four months,
what strikes me most now in these notes is awkwardness of the approach\xout{x}
followed at start,  when yielding\inserted{to} the reflex of lazyness of describing
abelianization of homotopy types via the semisimplicial grindmill.  The
uneasiness in these notes is obvious throughout  --  I kind of knew
``au fond'' that dragging in the particular test category \hand{$\Delta$} was rather
silly.  In section 100 only does it get clear that the best\hand{,}\xout{xxxx \unsure{intrinsic}}
description\hand{,}\inserted{for abelianization}within the modelizer (Cat),\deleted{for abelianization} is via the
unusual direct image $p_!^{ab}$ corresponding to the projection
\[ p=p_A \; : A \longrightarrow e \]
of the ``model'' A in (Cat) to the final object (formulae (11) and (12)
page 569),  by applying $p_!^{ab}$ to a ``cointegrator''  $L_{.}^{A}$ for A  namely to a
projective resolution of  $\underline{Z}_A$  (where A is written B by the way,  as \xout{we}
I had been led before to replace a given A in (Cat) by it's ``dual'' or
opposite B=A$^{o}$ ,  bound as \deleted{we were}\inserted{I was} for interpreting ``integrators'' for A
in terms of ``cointegrators'' for B ...);  and in the next section the
step is finally taken (against the ``block'' !)  to write
\[ p_!^{ab}(L_{.}^{B}) = Lp_!^{ab}(\underline{Z}_B) \quad , \]
which inserts abelianization into the familiar formalism of (left)
derived functors of standard functors.  The reasonable next thing to do
was of course \xout{to} what I finally did\inserted{only} yesterday,  namely check the commu-%
tativity of the diagram (7) of p. 563,  namely compatibility of this
notion of abelianization (or homology) with Whitehead's abelianization,
when working \xout{in} with models coming from A$^{\wedge}$,  A any object in (Cat).
This by the way,  when applied to the case  A = \hand{$\Delta$} ,  gives at once
the equivalence of the intrinsic definition of abelianization,  with
the semisimplicial one we started with  --  provided we remark that in
this case,  $Lp_!^{ab}$\inserted{(rather, it's restriction to $\hat{\Delta}_{ab}$)} may be equally interpreted as the Kan-Dold-Puppe
\margin{11}
